\begin{minipage}{0.74\linewidth}
À fin de vérifier la valeur de la concentration molaire $C_{A}$ de la solution
$\mathrm{(S_{A})}$, on réalise un titrage acido-basique.

Dans un bécher, on verse un volume $V_{A}=\qty{20,0}{\mL}$ de cette solution et
on y ajoute progressivement une solution aqueuse $\mathrm{(S_{B})}$ d'hydroxyde
de sodium $\ce{Na^+_{(aq)} + HO^-_{(aq)}}$ de concentration molaire
$C_{B}=\qty{0,25}{\M}$. Les coordonnées du point d'équivalence sont~:
$(V_{B,\mathrm{E}}=\qty{8,0}{\mL}~; \mathrm{pH_{E}}=8,2)$.

Le montage expérimental utilisé pour réaliser ce dosage est représenté sur
la figure ci-contre.
\end{minipage}
\hfill
\begin{minipage}{0.24\linewidth}
    \flushright
    \begin{tikzpicture}
        \node {\includegraphics[width=\textwidth]{2025_03_31_52b4902f5779608975f8g-2}};
    \end{tikzpicture}
\end{minipage}
\begin{enumerate}
    \item Nommer les éléments correspondants aux numéros indiqués sur le montage
          de la figure.
    \item Écrire l'équation de la réaction qui se produit entre l'acide
          méthanoïque $\ce{HCOOH_{(aq)}}$ et les ions hydroxydes
          $\ce{HO^-_{(aq)}}$ au cours du dosage, sachant qu'elle est totale.
    \item Vérifier la valeur de $C_{A}$.
    \item Parmi les deux indicateurs colorés suivants, quel est celui qui
          convient le mieux à ce dosage~? Justifier.

          \begin{minipage}{\linewidth}
              \begin{table}[H]
                  \centering
                  \renewcommand{\arraystretch}{1.5}
                  \begin{tabularx}{\textwidth}{|p{4cm}|C|C|C|}
                  \hline
                  \textbf{Indicateur coloré} & \textbf{Teinte acide} &
                  \textbf{Zone de virage} & \textbf{Teinte basique} \\
                  \hline
                  Rouge de crésol & Jaune & $7,2-8,8$ & Rouge \\
                  \hline
                  Alizarine & Rouge & $11,0-12,4$ & Violet \\
                  \hline
                  \end{tabularx}
              \end{table}
          \end{minipage}
    \item Pour un volume versé $V_{B}=\frac{V_{B,E}}{2}$ de la solution
          $\mathrm{(S_{B})}$, le $\mathrm{pH}$ du mélange dans le bécher vaut
          $\mathrm{pH}=3,8$ et $[\ce{HCOOH_{(aq)}}]=[\ce{HCOO^-_{(aq)}}]$.

          Calculer la constante d'acidité $\mathrm{K_{A}}$ du couple
          $(\ce{HCOOH_{(aq)}}\ttslash{}\ce{HCOO^-_{(aq)}})$.
\end{enumerate}


